\documentclass[a4paper,12pt]{article}

\usepackage[T2A]{fontenc}
\usepackage[utf8]{inputenc}
\usepackage[russian]{babel}
\usepackage{amsmath,amssymb,amsthm,mathtools}
\usepackage{helvet}
\renewcommand{\familydefault}{\sfdefault}
\usepackage{icomma}
\usepackage[a4paper,margin=20mm,headheight=15pt]{geometry}
\usepackage{microtype}
\usepackage{tikz}
\usetikzlibrary{calc,arrows.meta,positioning,shapes.geometric}
\usepackage{tkz-euclide}
\usepackage{pgfplots}
\pgfplotsset{compat=1.18}
\usepackage{tabularx,booktabs,longtable}
\usepackage{enumitem}
\usepackage{hyperref}
\usepackage{parskip}
\usepackage{fancyhdr}
\usepackage{setspace}
\usepackage{xcolor}

\definecolor{accent}{HTML}{008080}
\definecolor{darkaccent}{HTML}{004D4D}
\definecolor{textgray}{HTML}{333333}
\definecolor{softgray}{HTML}{F3F5F5}

\hypersetup{hidelinks}
\color{textgray}
\onehalfspacing
\setlength{\parindent}{0pt}
\setlist[itemize]{leftmargin=1.6em,itemsep=0.25em,topsep=0.3em}
\setlist[enumerate]{leftmargin=1.8em,itemsep=0.35em,topsep=0.3em}

\pagestyle{fancy}
\fancyhf{}
\fancyhead[L]{\small Чек-лист по математике}
\fancyhead[R]{\small София Кальней}
\fancyfoot[C]{\small \thepage}
\renewcommand{\headrulewidth}{0.3pt}

\newcommand{\sectionrule}{%
  \par
  \vspace{-0.35em}%
  {\color{accent}\rule{\linewidth}{0.6pt}}%
  \vspace{0.6em}%
}
\newcommand{\key}[1]{\textcolor{darkaccent}{\textbf{#1}}}

\tikzset{
  flowstart/.style={
    ellipse,
    draw=darkaccent,
    line width=0.9pt,
    align=center,
    minimum width=33mm,
    minimum height=10mm,
    fill=softgray
  },
  flowstep/.style={
    rounded corners=2mm,
    draw=darkaccent,
    line width=0.9pt,
    align=center,
    text width=50mm,
    minimum height=13mm,
    fill=white
  },
  flowdecision/.style={
    diamond,
    aspect=2.15,
    draw=darkaccent,
    line width=0.9pt,
    align=center,
    text width=38mm,
    inner sep=1.8pt,
    fill=softgray
  },
  flowarrow/.style={
    -{Latex[length=3mm]},
    line width=0.9pt,
    draw=darkaccent
  },
  flowlabel/.style={
    font=\small,
    fill=white,
    inner sep=1pt
  }
}

\begin{document}
\pagenumbering{gobble}

% =========================================================
% ТИТУЛЬНЫЙ ЛИСТ
% =========================================================

\begin{titlepage}
\thispagestyle{empty}

\begin{center}
\vspace*{22mm}

{\Large\textcolor{darkaccent}{\textbf{Чек-лист}}}\par
\vspace{5mm}

{\LARGE\textbf{Уравнения, текстовые задачи, графики и геометрия}}\par
\vspace{4mm}

{\large Подготовка к ОГЭ по математике}\par

\vspace{16mm}

{\large\textbf{София Кальней}}\par
\vspace{4mm}

{\normalsize Формат тренировочного блока: Путь А}\par

\vspace{14mm}

{\normalsize Дата: \today}\par

\vfill

{\normalsize Лёвин Артём Александрович эксклюзивно для Софии Кальней}\par

\vspace{9mm}

\begin{tikzpicture}[remember picture,overlay]
  \draw[accent!55,line width=1.2pt]
    ($(current page.south west)+(18mm,17mm)$)
      .. controls
      ($(current page.south west)+(65mm,29mm)$)
      and
      ($(current page.south east)+(-75mm,8mm)$)
      ..
      ($(current page.south east)+(-18mm,20mm)$);
\end{tikzpicture}

\end{center}
\end{titlepage}

\pagenumbering{arabic}

% =========================================================
% 1. УРАВНЕНИЯ
% =========================================================

\section*{1. Уравнения с чётными степенями и модулем}
\sectionrule

\key{Ключевая идея.}
Чётную степень удобно распознавать как квадрат другой степени:
\[
x^8=(x^4)^2,
\qquad
(20-x)^4=\bigl((20-x)^2\bigr)^2.
\]

Это позволяет применить разность квадратов или последовательно извлечь
арифметические квадратные корни.

\subsection*{Формулы, которые необходимо различать}

\[
\sqrt{a^2}=|a|
\qquad
\text{для любого } a\in\mathbb R,
\]

\[
(\sqrt a)^2=a
\qquad
\text{при } a\ge0.
\]

В первой формуле модуль обязателен: арифметический квадратный корень
всегда неотрицателен.

\subsection*{Пример: решить уравнение}

\[
x^8=(20-x)^4.
\]

\key{Способ 1: подгон под разность квадратов.}

\begin{align*}
x^8-(20-x)^4&=0,\\
\bigl(x^4-(20-x)^2\bigr)
\bigl(x^4+(20-x)^2\bigr)&=0.
\end{align*}

Второй множитель равен нулю только при одновременном выполнении
\[
x=0
\qquad\text{и}\qquad
20-x=0,
\]
что невозможно.

Следовательно,
\[
x^4-(20-x)^2=0.
\]

Ещё раз используем разность квадратов:
\[
\bigl(x^2-(20-x)\bigr)
\bigl(x^2+(20-x)\bigr)=0.
\]

Получаем
\[
x^2+x-20=0
\qquad\text{или}\qquad
x^2-x+20=0.
\]

Первое уравнение:
\[
x^2+x-20=(x+5)(x-4)=0,
\]
поэтому
\[
x=-5
\qquad\text{или}\qquad
x=4.
\]

Для второго уравнения
\[
D=(-1)^2-4\cdot1\cdot20=-79<0,
\]
поэтому действительных корней нет.

\key{Способ 2: через модуль.}

Обе части исходного уравнения неотрицательны, поэтому извлечение
арифметического квадратного корня даёт равносильное уравнение:
\[
x^4=(20-x)^2.
\]

Извлекаем корень ещё раз:
\[
x^2=|20-x|.
\]

Раскрываем модуль по определению:
\[
\begin{cases}
x^2=20-x, & x\le20,\\
x^2=x-20, & x>20.
\end{cases}
\]

В первой ветви:
\[
x^2+x-20=0,
\]
откуда
\[
x=-5,\qquad x=4.
\]

Оба значения удовлетворяют условию \(x\le20\).

Во второй ветви:
\[
x^2-x+20=0,
\]
и действительных корней нет.

\key{Ответ:}
\[
x\in\{-5,4\}.
\]

\subsection*{Контроль ошибок}

\begin{itemize}
  \item После \(\sqrt{A^2}\) записывайте \(|A|\), если знак выражения
  \(A\) заранее неизвестен.

  \item При раскрытии модуля обязательно фиксируйте условие каждой
  ветви и проверяйте найденные корни по этому условию.

  \item Сумма квадратов
  \[
  U^2+V^2
  \]
  равна нулю только при
  \[
  U=0
  \qquad\text{и}\qquad
  V=0
  \]
  одновременно.

  \item При разложении на множители следите за равносильностью
  преобразований.

  \item Не делите уравнение на выражение, которое может обращаться
  в ноль: так можно потерять корни.
\end{itemize}

% =========================================================
% БЛОК-СХЕМА 1
% =========================================================

\clearpage
\thispagestyle{fancy}

\begin{center}
{\Large\textcolor{darkaccent}{%
\textbf{Алгоритм: уравнение с чётными степенями и модулем}}}
\end{center}

\vspace{7mm}

\begin{center}
\begin{tikzpicture}[node distance=10mm and 14mm]

  \node[flowstart] (start) {Начало};

  \node[flowstep,below=of start] (rewrite) {
    Представьте чётные степени как квадраты:
    \[
    U^2=V^2
    \]
  };

  \node[flowdecision,below=12mm of rewrite] (choice) {
    Какой путь короче?
  };

  \node[
    flowstep,
    below left=16mm and 10mm of choice
  ] (factor) {
    Разность квадратов:
    \[
    (U-V)(U+V)=0
    \]
  };

  \node[
    flowstep,
    below right=16mm and 10mm of choice
  ] (root) {
    Извлеките корень и сведите задачу к модулю
  };

  \node[
    flowstep,
    below=17mm of factor
  ] (small1) {
    Решите полученные уравнения меньшей степени
  };

  \node[
    flowstep,
    below=17mm of root
  ] (cases) {
    Раскройте модуль по определению и решите каждую ветвь
  };

  \node[
    flowstep,
    below=21mm of choice,
    yshift=-55mm
  ] (check) {
    Проверьте условия ветвей, ОДЗ и отсутствие потери корней
  };

  \node[flowstart,below=of check] (finish) {
    Запишите ответ
  };

  \draw[flowarrow] (start) -- (rewrite);
  \draw[flowarrow] (rewrite) -- (choice);

  \draw[flowarrow]
    (choice)
    --
    node[flowlabel,above left]{факторизация}
    (factor);

  \draw[flowarrow]
    (choice)
    --
    node[flowlabel,above right]{корни и модуль}
    (root);

  \draw[flowarrow] (factor) -- (small1);
  \draw[flowarrow] (root) -- (cases);

  \draw[flowarrow]
    (small1.south)
    |-
    (check.west);

  \draw[flowarrow]
    (cases.south)
    |-
    (check.east);

  \draw[flowarrow] (check) -- (finish);

\end{tikzpicture}
\end{center}

\clearpage

% =========================================================
% 2. СМЕСИ И ВЫСУШИВАНИЕ
% =========================================================

\section*{2. Смеси, растворы и высушивание}
\sectionrule

\key{Основной принцип.}
В задачах на смеси сохраняется количество растворённого вещества.
В задачах на высушивание сохраняется количество сухого вещества.

\subsection*{Растворы}

Если концентрация дана в процентах, сначала переводите её
в десятичную дробь:
\[
26\%=0{,}26,
\qquad
44\%=0{,}44.
\]

Для стандартной математической модели с аддитивными объёмами:
\[
\text{количество вещества}
=
\text{объём раствора}
\cdot
\text{концентрация}.
\]

\begin{table}[h!]
\centering
\renewcommand{\arraystretch}{1.35}

\begin{tabularx}{\linewidth}{X c c c}
\toprule
&
Объём, л
&
Концентрация
&
Количество вещества
\\
\midrule

Первый раствор
&
5
&
\(0{,}26\)
&
\(5\cdot0{,}26\)
\\

Второй раствор
&
7
&
\(0{,}44\)
&
\(7\cdot0{,}44\)
\\

Смесь
&
12
&
\(x\)
&
\(12x\)
\\

\bottomrule
\end{tabularx}
\end{table}

Составляем уравнение по количеству растворённого вещества:
\[
12x
=
5\cdot0{,}26
+
7\cdot0{,}44.
\]

Вычисляем:
\[
12x=1{,}3+3{,}08=4{,}38.
\]

Следовательно,
\[
x=\frac{4{,}38}{12}=0{,}365.
\]

Переводим в проценты:
\[
0{,}365=36{,}5\%.
\]

\key{Ответ:}
\[
36{,}5\%.
\]

\textbf{Замечание.}
В реальной физике или химии при переходе между массой и объёмом
может потребоваться плотность. В стандартной математической задаче
используется модель, заданная условием.

\subsection*{Высушивание: виноград и изюм}

При высушивании уменьшается количество воды, а сухое вещество
сохраняется.

Если начальная масса равна \(M\), доля воды сначала равна \(p\),
а после высушивания --- \(q\), то сухое вещество составляет:

до высушивания
\[
M(1-p),
\]

после высушивания
\[
M_{\text{после}}(1-q).
\]

Поэтому
\[
M(1-p)
=
M_{\text{после}}(1-q).
\]

Отсюда
\[
M_{\text{после}}
=
\frac{M(1-p)}{1-q}.
\]

\subsection*{Контроль ошибок}

\begin{itemize}
  \item Не складывайте проценты концентраций напрямую.

  \item Перед составлением уравнения переводите проценты в доли:
  \[
  35\%=0{,}35.
  \]

  \item В задачах на высушивание определите, какая часть вещества
  сохраняется.

  \item Не путайте процент воды и процент сухого вещества:
  \[
  100\%-p\%
  \]
  составляет оставшаяся доля.

  \item В конце переводите десятичную дробь обратно в проценты,
  если этого требует вопрос.
\end{itemize}

% =========================================================
% 3. КУСОЧНЫЙ ГРАФИК
% =========================================================

\section*{3. Кусочно заданная функция и параметр \(m\)}
\sectionrule

Рассмотрим функцию
\[
f(x)=
\begin{cases}
2-|x|, & x\le4,\\
-x^2+10x-25, & x>4.
\end{cases}
\]

\subsection*{Первая ветвь: модуль}

Раскрываем модуль:
\[
2-|x|
=
\begin{cases}
2+x, & x\le0,\\
2-x, & x>0.
\end{cases}
\]

С учётом исходного ограничения \(x\le4\):
\[
f(x)=
\begin{cases}
2+x, & x\le0,\\
2-x, & 0<x\le4.
\end{cases}
\]

Для построения достаточно нескольких опорных точек:
\[
(-1,1),
\qquad
(0,2),
\qquad
(4,-2).
\]

Точка
\[
(4,-2)
\]
принадлежит графику, поскольку в первой ветви записано
\[
x\le4.
\]

\subsection*{Вторая ветвь: парабола}

Преобразуем:
\[
-x^2+10x-25
=
-(x^2-10x+25)
=
-(x-5)^2.
\]

Следовательно:

\[
x_{\text{вершины}}=5,
\qquad
y_{\text{вершины}}=0.
\]

Вершина:
\[
(5,0).
\]

Ветви параболы направлены вниз.

Вторая ветвь существует только при
\[
x>4.
\]

Если подставить \(x=4\) в полную параболу, получится:
\[
y=-(4-5)^2=-1.
\]

Точка
\[
(4,-1)
\]
выколота, поскольку значение \(x=4\) во вторую ветвь не входит.

\begin{center}
\begin{tikzpicture}
\begin{axis}[
  width=0.88\linewidth,
  height=8.2cm,
  axis lines=middle,
  xmin=-5.5,
  xmax=8.5,
  ymin=-5.5,
  ymax=3.2,
  samples=220,
  unit vector ratio=1 1,
  grid=both,
  minor tick num=1,
  clip=false,
  xlabel={$x$},
  ylabel={$y$},
  xtick={-4,-2,0,2,4,5,6,8},
  ytick={-4,-2,-1,0,1,2},
  every axis plot/.append style={line width=1.2pt}
]

\addplot[
  accent,
  domain=-5:0
]{x+2};

\addplot[
  accent,
  domain=0:4
]{2-x};

\addplot[
  darkaccent,
  domain=4.001:8.2
]{-(x-5)^2};

\addplot[
  only marks,
  mark=*,
  mark size=2.3pt,
  accent
]
coordinates {(4,-2)};

\addplot[
  only marks,
  mark=o,
  mark size=2.6pt,
  darkaccent,
  mark options={fill=white}
]
coordinates {(4,-1)};

\addplot[
  only marks,
  mark=*,
  mark size=2.3pt,
  darkaccent
]
coordinates {(5,0)};

\end{axis}
\end{tikzpicture}
\end{center}

\subsection*{Когда прямая \(y=m\) имеет ровно три общие точки?}

Прямая
\[
y=m
\]
является горизонтальной.

Для каждого положения такой прямой нужно подсчитать число
пересечений со всеми ветвями графика.

Особое внимание требуется пограничным уровням:
\[
m=-2,
\qquad
m=-1,
\qquad
m=0.
\]

Для данного графика ровно три общие точки возникают при
\[
m\in[-2,-1]\cup\{0\}.
\]

Почему границы нужно проверять отдельно:

\begin{itemize}
  \item при \(m=-2\) точка \((4,-2)\) включена;

  \item при \(m=-1\) потенциальная точка параболы с \(x=4\)
  выколота;

  \item при \(m=0\) горизонталь проходит через вершину
  параболы \((5,0)\).
\end{itemize}

% =========================================================
% БЛОК-СХЕМА 2
% =========================================================

\clearpage
\thispagestyle{fancy}

\begin{center}
{\Large\textcolor{darkaccent}{%
\textbf{Алгоритм: кусочный график и параметр \(m\)}}}
\end{center}

\vspace{7mm}

\begin{center}
\begin{tikzpicture}[node distance=6mm]

  \node[flowstart] (start) {
    Начало
  };

  \node[
    flowstep,
    below=of start
  ] (split) {
    Разделите функцию на ветви и выпишите ограничения на \(x\)
  };

  \node[
    flowstep,
    below=of split
  ] (module) {
    Раскройте модули и преобразуйте каждую формулу
    к удобному виду
  };

  \node[
    flowstep,
    below=of module
  ] (build) {
    Постройте каждую ветвь только на её области определения
  };

  \node[
    flowdecision,
    below=8mm of build
  ] (boundary) {
    Граница входит в ветвь?
  };

  \node[
    flowstep,
    below left=9mm and 10mm of boundary
  ] (filled) {
    Поставьте закрашенную точку
  };

  \node[
    flowstep,
    below right=9mm and 10mm of boundary
  ] (open) {
    Поставьте выколотую точку
  };

  \node[
    flowstep,
    below=14mm of boundary,
    yshift=-34mm
  ] (horizontal) {
    Проведите мысленно \(y=m\) и считайте общие точки
  };

  \node[
    flowstep,
    below=of horizontal
  ] (critical) {
    Отдельно проверьте уровни вершин и граничных точек
  };

  \node[
    flowstart,
    below=of critical
  ] (finish) {
    Запишите множество значений \(m\)
  };

  \draw[flowarrow] (start) -- (split);
  \draw[flowarrow] (split) -- (module);
  \draw[flowarrow] (module) -- (build);
  \draw[flowarrow] (build) -- (boundary);

  \draw[flowarrow]
    (boundary)
    --
    node[flowlabel,above left]{да}
    (filled);

  \draw[flowarrow]
    (boundary)
    --
    node[flowlabel,above right]{нет}
    (open);

  \draw[flowarrow]
    (filled.south)
    |-
    (horizontal.west);

  \draw[flowarrow]
    (open.south)
    |-
    (horizontal.east);

  \draw[flowarrow]
    (horizontal)
    --
    (critical);

  \draw[flowarrow]
    (critical)
    --
    (finish);

\end{tikzpicture}
\end{center}

\clearpage

% =========================================================
% 4. ГЕОМЕТРИЯ
% =========================================================

\section*{4. Геометрия: квадрат и окружность}
\sectionrule

Две соседние вершины квадрата \(A\) и \(B\) лежат на окружности,
а две другие вершины \(C\) и \(D\) лежат на её диаметре.
Радиус окружности равен \(4\).

\begin{center}
\begin{tikzpicture}[scale=0.78]

  \tkzSetUpLine[
    line width=1pt,
    color=accent
  ]

  \tkzSetUpPoint[
    color=accent,
    fill=accent,
    size=3
  ]

  \tkzSetUpLabel[
    color=black,
    font=\small
  ]

  \tkzDefPoint(0,0){O}
  \tkzDefPoint(4,0){R}
  \tkzDefPoint(-4,0){L}

  \tkzDefPoint(-1.7889,0){D}
  \tkzDefPoint(1.7889,0){C}
  \tkzDefPoint(-1.7889,3.5777){A}
  \tkzDefPoint(1.7889,3.5777){B}

  \tkzDrawCircle(O,R)
  \tkzDrawSegment[gray](L,R)

  \tkzDrawPolygon(A,B,C,D)
  \tkzDrawSegments(A,O B,O)

  \tkzDrawPoints(A,B,C,D,O)

  \tkzLabelPoints[above](A,B)
  \tkzLabelPoints[below](C,D,O)

  \tkzMarkRightAngle[size=0.25](A,D,O)
  \tkzMarkRightAngle[size=0.25](B,C,O)

  \tkzLabelSegment[above left](A,O){$4$}
  \tkzLabelSegment[left](A,D){$x$}
  \tkzLabelSegment[below](D,O){$\frac{x}{2}$}

\end{tikzpicture}
\end{center}

Проведём радиусы \(AO\) и \(BO\).

Треугольники \(ADO\) и \(BCO\) прямоугольные.

При этом
\[
AO=BO=4
\]
как радиусы окружности и
\[
AD=BC=x
\]
как стороны квадрата.

Следовательно, прямоугольные треугольники \(ADO\) и \(BCO\)
равны по гипотенузе и катету.

Поэтому
\[
DO=CO.
\]

Поскольку
\[
DC=x,
\]
получаем
\[
DO=CO=\frac{x}{2}.
\]

По теореме Пифагора в треугольнике \(ADO\):
\[
x^2+\left(\frac{x}{2}\right)^2=4^2.
\]

Искомая площадь квадрата:
\[
S=x^2.
\]

Поэтому удобно сразу найти \(x^2\):
\[
x^2+\frac{x^2}{4}=16,
\]
\[
\frac54x^2=16.
\]

Следовательно,
\[
x^2
=
16\cdot\frac45
=
\frac{64}{5}.
\]

Итак,
\[
S=\frac{64}{5}=12{,}8.
\]

\key{Ответ:}
\[
12{,}8.
\]

\subsection*{Полезная стратегия для задач с окружностью}

\begin{itemize}
  \item Начинайте чертёж с окружности и её центра.

  \item Если можно провести радиусы и получить прямоугольные
  треугольники, делайте это в первую очередь.

  \item Для двух треугольников сразу проверяйте, равны ли они,
  подобны ли они и какие равные элементы уже известны.

  \item Если требуется площадь \(x^2\), находите \(x^2\) сразу,
  не вычисляя \(x\) отдельно.
\end{itemize}

% =========================================================
% ТРЕНИРОВОЧНЫЙ БЛОК
% =========================================================

\clearpage

\section*{Тренировочный блок --- Путь А}
\sectionrule

\textbf{Решайте без подсказок.}
В графической задаче обязательно отмечайте включённые и выколотые
точки. В уравнениях проверяйте условия раскрытия модуля.

\begin{enumerate}

  \item
  Смешали \(4\) л \(15\%\)-го раствора и \(6\) л \(35\%\)-го
  раствора того же вещества.
  Найдите концентрацию получившегося раствора.

  \item
  В \(18\) кг винограда содержится \(90\%\) воды.
  После высушивания в изюме стало \(25\%\) воды.
  Найдите массу получившегося изюма.

  \item
  Решите уравнение
  \[
  x^8=(x+6)^4.
  \]

  \item
  Две соседние вершины квадрата лежат на окружности,
  две другие --- на её диаметре.
  Радиус окружности равен \(5\).
  Найдите площадь квадрата.

  \item
  Постройте график функции
  \[
  f(x)=
  \begin{cases}
  1-|x|, & x\le2,\\
  -(x-3)^2, & x>2.
  \end{cases}
  \]

  Найдите все значения \(m\), при которых прямая
  \[
  y=m
  \]
  имеет с графиком ровно три общие точки.

\end{enumerate}

% =========================================================
% ОТВЕТЫ
% =========================================================

\clearpage

\section*{Ответы}
\sectionrule

\begin{table}[h!]
\centering
\renewcommand{\arraystretch}{1.45}

\begin{tabularx}{0.82\linewidth}{c X}
\toprule
№ & Ответ \\
\midrule

1 & \(27\%\) \\

2 & \(2{,}4\) кг \\

3 & \(x\in\{-2,3\}\) \\

4 & \(20\) \\

5 & \(m\in\{-1,0\}\) \\

\bottomrule
\end{tabularx}
\end{table}

\vspace{8mm}

\textbf{Финальная самопроверка перед ОГЭ:}

\begin{itemize}

  \item переведены ли проценты в доли там, где это требуется;

  \item учтены ли ОДЗ и ограничения отдельных ветвей;

  \item сохранён ли модуль после
  \[
  \sqrt{A^2};
  \]

  \item проверены ли пограничные точки кусочного графика;

  \item исключено ли деление на ноль;

  \item не потеряны ли корни при сокращении;

  \item не появились ли лишние корни после возведения в квадрат;

  \item проверены ли знаки при переносах и раскрытии скобок;

  \item при работе с неравенствами изменён ли знак неравенства
  при умножении или делении на отрицательное число;

  \item правильно ли текст задачи переведён в уравнение;

  \item записан ли ответ в тех единицах и в том формате,
  которые требует условие.

\end{itemize}

\end{document}